\section{问题二：日前预测购电与日内滚动执行}

\subsection{问题分析}

问题二要求在\textbf{信息不完全}的条件下，逐日给出微网向外部电网的计划购电量，并在实际负荷、光伏与日前计划发生偏差时，用储能和紧急购电保证供电。与问题一的本质差别有三点。

第一，决策时刻被限制在每日 0:00。日期 \(D\) 的 144 个时段计划，只能使用 \(D\) 日 0:00 之前已经观测到的历史负荷与光伏，不得使用当天后续真值，也不使用附件 3。电价仍取附件 1 的典型日分时电价，且全年各日共用同一套 \(\pi_t\)；负荷与光伏的实际轨迹取自附件 2，仅在日内回放时打开。

第二，计划购电一旦形成即被锁死。实际净负荷高于计划时，不得按当天真值重解全天购电；缺口只能由储能放电或按当期电价 5 倍结算的紧急外购补足。紧急电只用于供电，不得用于充电。实际净负荷低于计划时，多余能量可给储能充电，充不进则弃光。微网仍不可向外网售电。

第三，评价窗口是 2025-02-01 至 12-31，共 334 天。2025-02-01 0:00 的实际储电量固定为 \(6000\,\mathrm{kWh}\)；1 月数据只用于预报训练和残差分位估计，不把 1 月滚动储电量带到 2 月 1 日。此后跨日传递的是回放得到的实际日末储电量，而不再强制问题一的日闭环 \(E_0=E_{24}=6000\,\mathrm{kWh}\)。计划轨迹与实际轨迹都必须落在运行窗 \([1200,10800]\,\mathrm{kWh}\) 内。

据此，问题二是一个两阶段决策：0:00 用预测净负荷做日前计划；日内在计划购电锁死的前提下，用实际净负荷滚动调度储能，并以计划费与 5 倍紧急费之和作为正式目标。5 倍罚价使欠计划的损失远大于多计划的损失，因此日前不宜把点预报当作真值，而应对负荷取上分位、对光伏取下分位，形成偏保守的净负荷。

\subsection{决策变量}

问题二将充、放电写在微网与储能的接口侧（交流侧）。接口充电 \(c_t\) 进入电池后变为 \(\eta c_t\)，电池放出 \(d_t/\eta\) 后在接口得到 \(d_t\)，其中 \(\eta=0.9\)。这与问题一的电池侧计量一一对应，只是符号落点不同。时段长度 \(\Delta t=1/6\,\mathrm{h}\)，\(t=1,\ldots,144\)。

日前阶段（0:00，仅用预测）的决策为
\[
G_t,\quad
c_t^{\mathrm{plan}},\quad
d_t^{\mathrm{plan}},\quad
w_t^{\mathrm{plan}},\quad
E_t^{\mathrm{plan}},
\]
其中 \(G_t\) 是向外部电网申报并锁死的计划购电量（kWh），\(c_t^{\mathrm{plan}}\)、\(d_t^{\mathrm{plan}}\) 是接口侧计划充、放电量（kWh），\(w_t^{\mathrm{plan}}\) 是计划弃光量（kWh），\(E_t^{\mathrm{plan}}\) 是时段末计划储电量（kWh）。

日内阶段（计划购电 \(G_t\) 已给定）的决策为
\[
c_t,\quad
d_t,\quad
w_t,\quad
G_t^{\mathrm{em}},\quad
E_t,
\]
其中 \(G_t^{\mathrm{em}}\) 是紧急购电量（kWh），\(E_t\) 是时段末实际储电量（kWh）。正式方案的日内层采用因果贪心规则，上述量为该规则的输出，而不是再解一次全天购电。

辅助量中，\(\hat{P}_{D,t}^{\mathrm{load}}\)、\(\hat{P}_{D,t}^{\mathrm{pv}}\) 为日期 \(D\) 时段 \(t\) 的点预报（kW）；经分位修正后记为 \(\hat{P}_{D,t}^{\mathrm{load},q}\)、\(\hat{P}_{D,t}^{\mathrm{pv},q}\)。实际功率记为 \(P_{D,t}^{\mathrm{load,act}}\)、\(P_{D,t}^{\mathrm{pv,act}}\)。电量由功率乘 \(\Delta t\) 得到。

\subsection{约束条件}

\noindent\textbf{（1）日前电量平衡与储电动态。}
对预测电量 \(\hat{L}_t=\hat{P}_{D,t}^{\mathrm{load},q}\Delta t\)、\(\hat{S}_t=\hat{P}_{D,t}^{\mathrm{pv},q}\Delta t\)，
\begin{equation}
\label{eq:q2-plan-balance}
G_t+\hat{S}_t+d_t^{\mathrm{plan}}
=\hat{L}_t+c_t^{\mathrm{plan}}+w_t^{\mathrm{plan}},
\end{equation}
\begin{equation}
\label{eq:q2-plan-soc}
E_t^{\mathrm{plan}}
=E_{t-1}^{\mathrm{plan}}
+\eta\,c_t^{\mathrm{plan}}
-\frac{d_t^{\mathrm{plan}}}{\eta},
\end{equation}
其中 \(E_0^{\mathrm{plan}}\) 取当日 0:00 已经实现的实际储电量。式~\eqref{eq:q2-plan-balance} 在预测口径下保证供电不低于预测负荷；日前不用紧急购电去“补”预测缺口。

\noindent\textbf{（2）容量、功率与符号约束。}
\begin{equation}
\label{eq:q2-bounds}
\begin{aligned}
1200&\le E_t^{\mathrm{plan}}\le 10800,\\
0&\le c_t^{\mathrm{plan}}\le P_{\max}\Delta t,\\
0&\le d_t^{\mathrm{plan}}\le P_{\max}\Delta t,\\
G_t&\ge 0,\qquad w_t^{\mathrm{plan}}\ge 0,
\end{aligned}
\end{equation}
其中 \(P_{\max}=5000\,\mathrm{kW}\) 为接口额定功率，故每时段接口电量上限为 \(P_{\max}\Delta t=833.33\,\mathrm{kWh}\)。\(\eta<1\) 时同时充放会凭空增加损耗，基线模型不加 0-1 互斥变量。

\noindent\textbf{（3）信息不集。}
\(\hat{P}_{D,t}^{\mathrm{load}}\)、\(\hat{P}_{D,t}^{\mathrm{pv}}\) 及其分位修正只依赖 \(\{d:d<D\}\) 的历史；\(G_t\) 一旦求出，日内不得改写。

\noindent\textbf{（4）日内实际平衡、锁死购电与紧急电。}
记实际负荷、光伏电量为 \(L_t^{\mathrm{act}}\)、\(S_t^{\mathrm{act}}\)，则
\begin{equation}
\label{eq:q2-actual-balance}
G_t+S_t^{\mathrm{act}}+d_t+G_t^{\mathrm{em}}
=L_t^{\mathrm{act}}+c_t+w_t,
\end{equation}
\begin{equation}
\label{eq:q2-actual-soc}
E_t=E_{t-1}+\eta c_t-\frac{d_t}{\eta},
\end{equation}
并且 \(c_t,d_t,w_t,G_t^{\mathrm{em}}\ge 0\)，接口功率与储电量上下限与~\eqref{eq:q2-bounds} 相同。紧急电不进入充电项。跨日衔接为 \(E_0(D+1)=E_{144}(D)\)（实际值）；仅在 \(D=\)2025-02-01 时令 \(E_0=6000\,\mathrm{kWh}\)。

\noindent\textbf{（5）因果贪心可行域。}
正式日内规则在~\eqref{eq:q2-actual-balance}--\eqref{eq:q2-actual-soc} 上再加当期信息约束：第 \(t\) 个时段只根据当期实际净负荷与当期储电量决定 \(c_t,d_t,G_t^{\mathrm{em}},w_t\)，不预支未来真值。记残差 \(R_t=L_t^{\mathrm{act}}-S_t^{\mathrm{act}}-G_t\)。若 \(R_t\le 0\)，则按接口与容量上限吸收盈余，剩余弃光，且 \(d_t=G_t^{\mathrm{em}}=0\)；若 \(R_t>0\)，则先放电至功率与下限允许的最大值，不足部分记为 \(G_t^{\mathrm{em}}\)，且 \(c_t=0\)。该规则自动满足供电约束，且不同时充放。

\subsection{目标函数}

日前层在预测电量上最小化计划购电费：
\begin{equation}
\label{eq:q2-obj-plan}
\min\sum_{t=1}^{144}\pi_t G_t.
\end{equation}
式~\eqref{eq:q2-obj-plan} 不含日末储电量奖励。试算表明，若添加 \(-\mu E_{144}^{\mathrm{plan}}\) 且 \(\mu\eta<\min_t\pi_t\)，则谷电套利动机不足，日前计划与 \(\mu=0\) 重合；若 \(\mu\) 过大，则会在低谷把储电量顶到 \(10800\,\mathrm{kWh}\)。因此正式方案取 \(\mu=0\)。

日内层不重解 \(G_t\)。实际结算费用为
\begin{equation}
\label{eq:q2-obj-settle}
J_D=\sum_{t=1}^{144}\pi_t G_t+\sum_{t=1}^{144}5\pi_t G_t^{\mathrm{em}}.
\end{equation}
正式评价指标是 2025-02-01 至 12-31 的 \(\sum_D J_D\)。点预报误差会同时抬高~\eqref{eq:q2-obj-plan} 的计划费和~\eqref{eq:q2-obj-settle} 的紧急费，但紧急费单位成本约为计划费的 5 倍，故选预报口径时以总费用为准，不以 MAE/RMSE 为准。

\subsection{模型求解}

\begin{enumerate}[label=\textbf{步骤\arabic*},leftmargin=2.6em,itemsep=0.45em]
\item \textbf{数据对齐与门禁。}
将附件 1、附件 2 按时段终点对齐为全年 \(365\times 144\) 条记录，功率换算为电量 \(P\Delta t\)。检查时间连续、电价与问题一共用、负荷为正、光伏非负。1 月样本保留为预报历史，正式费用从 2 月 1 日起算。

\item \textbf{因果点预报。}
对日期 \(D\) 的每个时刻 \(t\)，用 \(d<D\) 的历史构造特征（时刻、星期、月份、年积日，以及同刻滞后 1/2/7 日与 3/7/30 日滑动统计）。负荷采用 Expanding 窗口的 XGBoost 回归（\(n_{\mathrm{estimators}}=200\)，最大深度 6，学习率 0.05，随机种子 42）；光伏改用过去 7 日同刻均值。预报截断为非负，夜间历史最大光伏接近 0 的时刻强制为 0。对照试验中，光伏 XGB 的点预测 MAE 并不优于 7 日均值，故正式骨干取“负荷 XGB + 光伏 7 日均值”。

\item \textbf{残差分位修正。}
对已过去的日期，计算同刻残差 \(r_{d,t}^{L}=P_{d,t}^{\mathrm{load,act}}-\hat{P}_{d,t}^{\mathrm{load}}\)、\(r_{d,t}^{\mathrm{PV}}=P_{d,t}^{\mathrm{pv,act}}-\hat{P}_{d,t}^{\mathrm{pv}}\)，取
\[
\hat{P}_{D,t}^{\mathrm{load},q}
=\max\bigl(0,\hat{P}_{D,t}^{\mathrm{load}}+Q_t^{L}(0.8)\bigr),\qquad
\hat{P}_{D,t}^{\mathrm{pv},q}
=\max\bigl(0,\hat{P}_{D,t}^{\mathrm{pv}}+Q_t^{\mathrm{PV}}(0.2)\bigr),
\]
其中 \(Q_t^{L}(0.8)\)、\(Q_t^{\mathrm{PV}}(0.2)\) 为 \(d<D\) 同刻残差的经验分位。负荷上分位、光伏下分位使日前净负荷偏高，从而多买计划电、少触发 5 倍紧急电。该修正会恶化点预测 MAE，这是费用非对称下的有意取舍。

\item \textbf{求解日前线性规划。}
将修正后的 144 点电量、当日 0:00 实际储电量与电价 \(\pi_t\) 代入约束~\eqref{eq:q2-plan-balance}--\eqref{eq:q2-bounds} 和目标~\eqref{eq:q2-obj-plan}，用 CBC 求最优解，得到锁死的 \(\{G_t\}\)。求解状态须为 Optimal；计划轨迹须满足电量平衡、SOC 窗和接口功率上限。

\item \textbf{锁死 \(G_t\) 后的日内贪心回放。}
按时段打开附件 2 的当期真值，用步骤 4 的 \(G_t\) 计算 \(R_t\)，执行上文因果贪心规则，得到 \(c_t,d_t,w_t,G_t^{\mathrm{em}},E_t\)，并按~\eqref{eq:q2-obj-settle} 累加当日费用。回放不得使用未来时段的实际负荷或光伏。对照中还实现了剩余时段模型预测控制：当期用实测、其余时段用日前预报，重解剩余储能与紧急电的线性规划，但只执行第一段。该规则在阶段 1 上未能降低总费用，正式方案不采用。

\item \textbf{跨日传递与正式窗口。}
2025-02-01 0:00 将实际储电量重置为 \(6000\,\mathrm{kWh}\)；此后 \(E_{144}\) 作为次日 \(E_0\)。1 月不进入 SOC 仿真。正式统计只保留 2 月 1 日至 12 月 31 日。

\item \textbf{阶段 1 消融选策。}
在 2025-02-01 至 02-14 上比较点预报贪心（V1）、日末留电奖励、剩余时段 MPC、以及分位保守预报。先知下界（用当天真值做日前 LP，不作正式成绩）为 \(585738.41\) 元。结果见表~\ref{tab:q2-phase1}。\(\mu\in\{0.25,0.4\}\) 与 V1 完全重合；仅改 MPC 时总费用略降但紧急时段增多；分位方案 \(\mathrm{C}\) 把紧急费从 \(64782.27\) 元降到 \(182.27\) 元，总费用最低，因而定为正式策略。

\item \textbf{全年回放、校验与填表。}
用正式策略回放 334 天。校验包括：日前与回放电量平衡、SOC 落在 \([1200,10800]\)、接口功率不超过 \(5000\,\mathrm{kW}\)、\(G_t\ge 0\)、不同时充放、无未供电量、2 月 1 日实际初值恰为 \(6000\,\mathrm{kWh}\)。全部通过后写入 \(\texttt{result2.xlsx}\)：计划购电为日前锁死的 \(G_t\)（时间段改写为物理区间）；充放电与 0:00/24:00 储电量为回放实际值（交流侧）；紧急购电把相邻 10 分钟槽合并为连续时段。
\end{enumerate}

阶段 1 与全年的冻结数字分别见表~\ref{tab:q2-phase1}、表~\ref{tab:q2-year}。正式策略全年总费用为 \(13631166.17\) 元，其中计划费 \(13211475.24\) 元、紧急费 \(419690.92\) 元（\(71414.67\,\mathrm{kWh}\)）。同一代码、同一 SOC 口径下的点预报对照为 \(14830254.38\) 元，总费用下降约 \(8.1\%\)，紧急费由 \(2612040.11\) 元降至 \(419690.92\) 元。计划日末储电量仍常落在下限 \(1200\,\mathrm{kWh}\) 附近，实际日末均值约 \(1915\,\mathrm{kWh}\)，说明费用下降主要来自分位保守带来的计划电前置，而不是日闭环或留电奖励。校验错误数为 0。

题面指定日 2025-03-20、06-21、09-23、12-21 的计划购电见表~\ref{tab:q2-specified}。四日中仅 9 月 23 日 20:20--20:30 出现紧急购电 \(6.38\,\mathrm{kWh}\)，其余三日紧急费为 0。完整 144 点策略见 \(\texttt{result2.xlsx}\)。

\begin{table}[htbp]
  \centering
  \caption{问题二阶段 1（2025-02-01 至 02-14）消融结果。先知下界不作正式成绩。}
  \label{tab:q2-phase1}
  \begin{tabular}{lrrrr}
    \toprule
    策略 & 总费用/元 & 计划费/元 & 紧急费/元 & 2/1 SOC \\
    \midrule
    点预报+贪心（V1） & 658945.08 & 594162.82 & 64782.27 & 6000 \\
    留电 \(\mu=0.25/0.4\) & 658945.08 & 594162.82 & 64782.27 & 6000 \\
    点预报+MPC & 654043.03 & 594072.45 & 59970.59 & 6000 \\
    分位保守+贪心（正式） & 646140.91 & 645958.64 & 182.27 & 6000 \\
    分位保守+MPC & 647583.53 & 645943.23 & 1640.29 & 6000 \\
    先知下界 & 585738.41 & 585738.40 & 0.01 & 6000 \\
    \bottomrule
  \end{tabular}
\end{table}

\begin{table}[htbp]
  \centering
  \caption{问题二正式窗口 2025-02-01 至 12-31（334 天）费用对照。对照方案与正式方案均从 \(6000\,\mathrm{kWh}\) 起算。}
  \label{tab:q2-year}
  \begin{tabular}{lrrrrr}
    \toprule
    方案 & 总费用/元 & 计划费/元 & 紧急费/元 & 紧急电量/kWh & 2/1 SOC \\
    \midrule
    点预报对照 & 14830254.38 & 12218214.27 & 2612040.11 & 434246.77 & 6000 \\
    正式分位方案 & 13631166.17 & 13211475.24 & 419690.92 & 71414.67 & 6000 \\
    \bottomrule
  \end{tabular}
\end{table}

\begin{table}[htbp]
  \centering
  \caption{指定日计划购电量（kWh）及当日费用（元）。购电量为日前锁死值。}
  \label{tab:q2-specified}
  \begin{tabular}{lrrrr}
    \toprule
    时段 & 03-20 & 06-21 & 09-23 & 12-21 \\
    \midrule
    10:00--10:10 & 0.00 & 0.00 & 0.00 & 0.00 \\
    12:00--12:10 & 645.50 & 0.00 & 498.55 & 1035.47 \\
    14:00--14:10 & 0.00 & 0.00 & 0.00 & 0.00 \\
    16:00--16:10 & 520.16 & 184.04 & 553.30 & 888.49 \\
    18:00--18:10 & 710.64 & 0.00 & 780.16 & 727.60 \\
    20:00--20:10 & 0.00 & 0.00 & 0.00 & 0.00 \\
    \midrule
    全日计划购电/kWh & 69928.90 & 34724.71 & 70395.52 & 99003.00 \\
    计划费/元 & 42547.03 & 19982.57 & 43819.10 & 64032.73 \\
    紧急费/元 & 0.00 & 0.00 & 43.00 & 0.00 \\
    总费用/元 & 42547.03 & 19982.57 & 43862.10 & 64032.73 \\
    \bottomrule
  \end{tabular}
\end{table}
